% Generated draft: requires mathematical and source review.
\documentclass[11pt]{article}
\usepackage[margin=25mm]{geometry}
\usepackage{fontspec}
\setmainfont{Latin Modern Roman}
\usepackage{amsmath,amssymb,mathtools,unicode-math}
\setmathfont{Latin Modern Math}
\usepackage{enumitem,needspace,xurl,hyperref}
\setlength{\parindent}{0pt}
\setlength{\parskip}{5pt}
\setlength{\emergencystretch}{3em}
\newcommand{\N}{\mathbb{N}}
\newcommand{\Z}{\mathbb{Z}}
\newcommand{\Q}{\mathbb{Q}}
\newcommand{\R}{\mathbb{R}}
\newcommand{\C}{\mathbb{C}}
\newcommand{\F}{\mathbb{F}}
\newcommand{\A}{\mathbb{A}}
\newcommand{\PP}{\mathbb{P}}
\newcommand{\E}{\mathbb{E}}
\newcommand{\eps}{\varepsilon}
\newcommand{\dd}{\,\mathrm{d}}
\begin{document}
\section*{UnsolvedMath ID: 800862}
\textbf{Source problem code:} catalogue entry
\subsection{Definitions and mathematical statement}
Let \(K\) denote a field, which the supplied source does not specify, and let \(I=\{1,2,3,4\}\). A \(4\times4\times4\) table is an array \(p=(p_{ijk})_{i,j,k\in I}\in K^{64}\), regarded as a tensor in \(K^4\otimes_K K^4\otimes_K K^4\). Its tensor rank over \(K\) is the least nonnegative integer \(r\) for which there exist vectors \(a^{(s)},b^{(s)},c^{(s)}\in K^4\), for \(1\le s\le r\), such that
\[
p_{ijk}=\sum_{s=1}^{r}a_i^{(s)}b_j^{(s)}c_k^{(s)}
\qquad(i,j,k\in I).
\]
Define
\[
S_K=\{p\in K^{64}:\operatorname{rank}_K(p)\le4\},
\qquad V_K=\overline{S_K}^{\mathrm{Zar}}.
\]
The closure is introduced explicitly because the source calls the rank-at-most-four locus a variety, without specifying whether it means the rank locus itself or its Zariski closure. These meanings cannot be silently identified.

Let \(\mathcal P_K\subseteq K[p_{ijk}:i,j,k\in I]\) stand for the particular collection of known polynomials of degree at most nine mentioned in the source. The source asserts that these polynomials vanish on its intended variety, but the supplied record neither lists nor characterizes this collection.

\textbf{Question, conditional on the closure interpretation.} For the intended field \(K\) and the intended collection \(\mathcal P_K\), is
\[
V_K=\{p\in K^{64}:f(p)=0\text{ for every }f\in\mathcal P_K\}?
\]
This displays the set-theoretic interpretation of ``cut out.'' If the source instead intends the rank locus itself, the left side must be \(S_K\). If it intends scheme-theoretic defining equations, equality of zero sets does not express that stronger target. The field, the polynomial collection, and these interpretations remain unresolved.

Although the source describes contingency tables and suggests DNA bases as index labels, it imposes no explicit nonnegativity or normalization condition. No such condition has been added here.
\subsection{Short English statement}
For \(4\times4\times4\) tensors of rank at most four, do the particular known equations of degree at most nine mentioned by Sturmfels define the intended variety? The supplied record leaves the field, the equations, and the precise meaning of the variety and of ``cut out'' unspecified.
\subsection{Sources}
\begin{itemize}
\item Supplied local UnsolvedMath registry record, permanent ID \(800862\), exact source problem code \texttt{catalogue entry}, titled \emph{Open Problems in Algebraic Statistics --- Specific Problem 1}; collection: AMR Open Problem Lists. The record attributes the problem to Sturmfels, with proposed year 2007, and labels its status as open. This is a supplied status label, not an independently verified literature assessment.
\item Sturmfels, \emph{Open Problems in Algebraic Statistics} (2007), Specific Problem (\#1). Supplied source link: \url{https://arxiv.org/abs/0707.4558}. The background reports access on 2026-10-16 and extraction by \texttt{arxiv\_latex}; these provenance details are supplied and have not been independently verified. The extracted statement contains a stray closing brace, treated as a formatting artifact.
\end{itemize}
\end{document}
